\documentclass[a4paper,10pt]{article}
\usepackage[left=3cm,top=3cm,right=2.5cm,bottom=4cm]{geometry}
\usepackage[utf8]{inputenc}
\usepackage[activeacute,spanish]{babel}
\usepackage{multirow}	% Para poder unir filas en las tablas
\usepackage{colortbl}	% Para colorear tablas
\usepackage{latexsym,amssymb,oldlfont,amsmath}
\usepackage{graphicx,graphics}
\usepackage{multicol}
\usepackage{color}
\usepackage{listings}
\usepackage{lastpage}
\usepackage{enumitem}



%\usepackage{fontspec}
%\setmainfont{Times New Roman}
\newcommand{\R}{{\mathbb R}}
\newcommand{\si}{{\mathcal S}^n}
\renewcommand{\theequation}{\thesection.\arabic{equation}} % Cambia formato salida numeracion equation
\makeatletter\@addtoreset {equation}{section} \makeatother % Tambien se usa
\renewcommand{\baselinestretch}{1.3}
\usepackage{amssymb}
\usepackage{hyperref}
\definecolor{midnightblue}{rgb}{0.1,0.1,0.44}


\newcommand{\dis}{\displaystyle}
\renewcommand{\theenumii}{\arabic{enumii}}

\def\Z{\mathbb{Z}}
\def\N{\mathbb{N}}
\def\R{\mathbb{R}}
\def\M{\mathcal{M}_n(\mathbb{R})}
\def\sin{\operatorname{sen}}
\def\e{\operatorname{e}}
\DeclareMathOperator{\arcsec}{arcsec}
\DeclareMathOperator{\arccot}{arccot}
\DeclareMathOperator{\arccsc}{arccsc}



\usepackage{etoolbox,fancyhdr,xcolor}
\fancyhf{} % sets both header and footer to nothing
\renewcommand{\headrulewidth}{0pt}
%\renewcommand{\footrulewidth}{0.5pt}
\pagestyle{fancy}
\headheight=60pt %para cambiar el tamaño del encabezadox
\fancyhead[L] %la "L" indica a la izquierda
{
\begin{minipage}{6cm}
\includegraphics[width=0.9\textwidth]{logo_udp.eps}
\end{minipage}
}

%\fancyhead[C]
%{
%\textbf{Cálculo II - Guía 1}
%}


%\fancyhead[R] %la "L" indica a la izquierda
%{

%}
%\begin{minipage}{4cm}
%\small %tamaño de letra pequeño
%{
%\textsf
%{
%Universidad de Chile\\
%...}}
%\end{minipage}}



%\fancyfoot[C]{}
%{
%\begin{minipage}{10cm}
%\color{midnightblue}
%\begin{center}Casa Matriz: Monjitas 527 Piso 15 - Santiago de Chile\\
%Fono (02)24367255 \end{center}
%\end{minipage}
%}
\fancyfoot[C]{}
%Página \thepage\ de \pageref{LastPage}}% Páginas Pares numeradas a la derecha
%\fancyfoot[LO]{\thepage} % Páginas Impares numeradas a la izquierda

\title{C2_Ay4}

\begin{document}

\begin{center}
Cálculo II\\
Ayudantía Nº4\\
Primer Semestre 2017
\end{center}



\begin{enumerate}
\item Justifique la veracidad del siguiente enunciado:

Una empresa tiene como función de ingreso marginal a $I'(x)=35-4x$ y función de costo marginal a $C'(x) = 5-2x$, donde $x$ es la cantidad de artículos vendidos, y las unidades monetarias están en millones de pesos [MM\$]. Suponga que cuando no se venden artículos el ingreso total es cero y el costo total es $25$[MM\$] (por conceptos de gastos fijos). Recordando que la utilidad de una empresa es la diferencia entre los ingresos totales y los costos totales, entonces la cantidad de artículos a vender que maximizan las utilidades son $225$ artículos.

\item Considere la función $f(x) = |\sin(x)|$ definida para $x \in [0,2\pi]$, y la partición $\displaystyle P = \left\{ 0, \frac{\pi}{4} , \frac{2\pi}{3} , \pi , \frac{3\pi}{2} , \frac{11\pi}{6} , 2\pi \right\}$. Calcule la suma superior e inferior de $f$ sobre $P$.


\item Considere la función $f(x)= [\sqrt{x}]$, donde $[x]$ denota la parte entera de $x$.
\begin{enumerate}[label=\emph{\alph*)}]
\item Grafique la función para $x \in [0,16]$
\item Obtenga una partición regular P del intervalo $[0,16]$ en 8 subintervalos.
\item  Calcule la suma superior e inferior de $f$ acotada a la partición P.
\end{enumerate}

\item Mediante sumas de Riemann, calcule el área bajo la curva $f(x) = 9-x^2$ en el intervalo $[-3,3]$

\item Mediante sumas de Riemann, calcule el área bajo la curva $f(x) = x^3$ en el intervalo $[0,1/2]$

\end{enumerate}




\vspace{10mm}




\end{document}
